\documentclass{article}
\usepackage{lmodern}

\usepackage{graphicx} % Required for inserting images
\usepackage{tcolorbox}

\usepackage[ruled,vlined]{algorithm2e}
\usepackage{booktabs}
\usepackage{amsmath}
\usepackage{mathtools}
\usepackage{amssymb}
\usepackage{tabularx}
\usepackage{float}

\usepackage[round,authoryear]{natbib}
\usepackage{amsthm}

\newtheorem{definition}{Definition}
\newtheorem{proposition}{Proposition}
\newtheorem{theorem}{Theorem}
\newtheorem{lemma}{Lemma}
\newtheorem{remark}{Remark}
\newtheorem{corollary}{Corollary}

\usepackage[
    top=1in,
    bottom=1in,
    left=1in,
    right=1in,
    includehead,
    includefoot
]{geometry}

\usepackage[hidelinks]{hyperref} % keep last


\renewcommand{\algorithmautorefname}{Algorithm}


\title{Scalable and Training-Free Ranking from Pairwise Comparisons via Acyclic Graph Construction}



\author{Soroush Vahidi\\
New Jersey Institute of Technology, Newark, NJ, USA\\
\texttt{sv96@njit.edu}\\
\href{https://orcid.org/0000-0003-1934-6282}{ORCID: 0000-0003-1934-6282}}




\date{}

\begin{document}



\maketitle

\begin{abstract}
Deriving global rankings from noisy and incomplete pairwise comparisons is a fundamental problem with applications in sports analytics, preference aggregation, and evaluation. Many recent approaches rely on learning-based models, but in practice ranking systems are often required to operate under strict constraints on training cost, runtime, and reproducibility. We study a scalable, training-free alternative that constructs rankings directly from weighted directed comparison graphs.

Our approach represents pairwise comparisons as a weighted digraph and leverages the connection between ranking inconsistency and feedback-arc-set removal to build an acyclic comparison backbone. The proposed pipeline is time-bounded and consists of three stages: a local-ratio-style cycle-breaking heuristic, a stable weight-prioritized add-back procedure with up to three passes while preserving acyclicity, and an optional bounded score-refinement stage under upset-based objectives used in recent ranking benchmarks. The resulting method outputs a real-valued score vector whose induced order is consistent with the recovered acyclic structure.

Across the benchmark suite used in recent ranking-from-comparisons work, the proposed method is competitive with strong classical baselines and delivers substantial runtime advantages over training-based GNNRank configurations under fixed wall-clock budgets. These results position the method as a practical, deterministic, and scalable alternative for ranking on large pairwise-comparison graphs when training cost and deployment efficiency are important.
\end{abstract}
\newcommand{\keywords}[1]{\textbf{Keywords:} #1}

\keywords{Ranking; Pairwise comparisons; Feedback arc set; Directed graphs; Training-free methods}
\section{Introduction}
\label{sec:intro}
Ranking from pairwise comparisons arises throughout statistics, machine learning, and network analysis, from sports analytics to preference aggregation. Methods typically return a real-valued score vector whose induced order is evaluated by upset-based losses \citet{he22}. Classical estimators are fast and training-free; learning-based directed GNNs such as GNNRank can be accurate but introduce training cost, hyperparameter sensitivity, and reduced reproducibility.

Directed comparison graphs almost always contain cycles arising from noise, incompleteness, and local inconsistencies. A combinatorial response is to delete a set of directed arcs so that the remaining graph is acyclic, then read a ranking from a topological order of the resulting directed acyclic graph (DAG). This is the Minimum Weighted Feedback Arc Set (MWFAS) viewpoint: minimize the total weight of deleted comparisons subject to acyclicity. Related spectral and synchronization methods already note connections between ranking inconsistency and feedback structure \citet{SerialRank,SyncRank}. MWFAS is NP-hard \citet{K72}, so exact minimization is generally impractical, yet the formulation supplies a precise consistency-restoration objective for ranking.

We emphasize at the outset that this paper does \emph{not} propose a new general approximation algorithm for MWFAS. Our objective is a practical, training-free, deterministic, and time-bounded ranking pipeline for weighted pairwise-comparison graphs. Prior work by \citet{VK25} already operationalized ranking via the local-ratio cycle-breaking heuristic of Demetrescu and Finocchi \citet{CI03}, together with weight-prioritized reinsertion of deleted arcs under an exact cycle-safety test and optional order-preserving score refinement. The present manuscript clarifies that lineage, corrects an implementation-fidelity gap, adds theory and a secondary structural repair, and expands empirical validation.

To make the ranking--MWFAS connection precise, we begin with formal definitions.

\begin{definition}[Ranking Problem]
Let $G=(V,E)$ be a directed graph, where $V$ is the set of items and $E$ is the set of directed pairwise-comparison edges.
Each edge $(i,j)\in E$ has weight $w_{ij}>0$, representing the strength of the preference for $i$ over $j$.
A \emph{ranking} is a total order on $V$, represented equivalently by a bijection $R:V\rightarrow\{1,\dots,|V|\}$, where smaller values indicate higher rank.
An edge $(i,j)$ is \emph{inconsistent} with $R$ if $R(i)>R(j)$.
The goal is to find a ranking that minimizes the total weight of inconsistent edges:
\begin{equation}
\label{eq:ranking_obj}
\min_{R}\ \sum_{(i,j)\in E:\ R(i)>R(j)} w_{ij}.
\end{equation}
\end{definition}

\begin{definition}[Minimum Weighted Feedback Arc Set (MWFAS)]
Let $G=(V,E)$ be a directed graph with edge weights $w_{ij}>0$ for each $(i,j)\in E$.
A set of edges $F\subseteq E$ is a \emph{feedback arc set} if removing $F$ makes the graph acyclic, i.e., $(V,E\setminus F)$ is a directed acyclic graph (DAG).
The MWFAS problem is
\begin{equation}
\label{eq:mwfas}
\min_{F\subseteq E}\ \sum_{(i,j)\in F} w_{ij}
\quad \text{s.t.} \quad (V,E\setminus F)\ \text{is acyclic.}
\end{equation}
\end{definition}

The ranking objective~\eqref{eq:ranking_obj} and MWFAS~\eqref{eq:mwfas} share the same optimum \emph{value}: the backward edges of any ranking form a feedback arc set of equal weight, and any feedback arc set induces a DAG whose topological orders realize that deleted weight as ranking disagreement. The correspondence between optimal rankings and optimal feedback arc sets is many-to-many rather than one-to-one; we treat this equivalence as a formal claim to be stated and proved, not merely asserted.

\subsection{Scope, lineage, and contributions}
The pipeline studied here (\S\ref{sec:framework}) builds an acyclic backbone with the local-ratio heuristic of \citet{CI03}, reinserts high-weight arcs while preserving acyclicity, extracts scores from the resulting DAG, and optionally applies order-preserving refinement. Several components are \emph{prior}: local-ratio Phase~A, exact cycle-safe add-back (including weight ordering), and ternary/refinement ideas already appear in \citet{CI03,VK25}. What this paper contributes is therefore not a reinvention of that core, but the following.

\begin{enumerate}
\item \textbf{Formalization and guarantee boundaries.}
We prove ranking--MWFAS optimum-value equivalence under the many-to-many correspondence above, restate implementation complexity carefully, and audit when the Demetrescu--Finocchi guarantee does and does not transfer. In particular, a practical timeout-bounded run does \emph{not} automatically inherit a nontrivial general approximation guarantee under early termination.
\item \textbf{Implementation-fidelity correction.}
The submitted codebase used a weaker fixed-topological-position proxy in place of the intended exact reachability test for cycle-safe reinsertion. Restoring exact reachability is fidelity to the established specification, not a newly invented mechanism; controlled ablations show that the correction materially improves ranking quality relative to the proxy.
\item \textbf{Secondary structural repair.}
We introduce a conflict-triggered weighted min-cut exchange with proofs of acyclicity preservation, strict improvement of the weighted feedback objective, and finite termination under strict improvement. Empirically, its benefit is regime-specific rather than universal.
\item \textbf{Expanded empirical validation.}
We compare against classical and GNNRank baselines under a common protocol with timeout-aware coverage, paired statistics, family-aware sensitivity (addressing Basketball-dominated suites), and a full ablation/sensitivity/scalability campaign. The modest expansion of the dataset suite relative to prior work is not presented as a primary novelty.
\end{enumerate}

Empirically, ranking quality shows strong gains against several classical and GNN methods, while remaining competitive---not uniformly superior---with \textsc{SpringRank}, \textsc{DavidScore}, and \textsc{SVD\_NRS} under family-aware analysis; \textsc{BTL} remains stronger on \texttt{upset\_ratio}. On runtime, the pipeline is slower than lightweight classical baselines and substantially faster than trained GNNRank models under the evaluated end-to-end training protocol. Scalability is strong on the evaluated sparse and moderately dense instances; the near-complete Finance graph is an explicit large-dense boundary case.

\subsection{Related Work}
\label{sec:related}

\subsubsection{Feedback Arc Set and MWFAS}
The Minimum Weighted Feedback Arc Set (MWFAS) problem asks for a minimum-weight set of directed arcs whose removal makes a graph acyclic \citet{K72,FPP20}. Classical approximation results include divide-and-conquer methods based on spreading metrics and LP relaxations \citet{EG95,EG2000}. Exact and parameterized algorithms target smaller or structurally restricted instances \citet{BA21,HM21,RS06,SJ21}. Heuristic refinements emphasize minimal/stable feedback sets and initialization effects \citet{c24,c25}. Closely related is the maximum acyclic subgraph problem \citet{HR94}. We use MWFAS as a ranking consistency objective, not as a setting in which we claim a new general approximation algorithm.

\subsubsection{Ranking from Pairwise Comparisons}
Ranking from noisy and incomplete pairwise comparisons has been studied through statistical, spectral, random-walk, and optimization-based formulations \citet{H00,CGZ21,SW16}. Our experiments evaluate score-based outputs under the upset-style losses used in recent ranking benchmarks \citet{he22}, so we focus on the principal baseline families represented in that protocol.

\subsubsection{Statistical and Spectral Ranking}
\textsc{BTL} models pairwise win probabilities through latent abilities \citet{BTL}. \textsc{DavidScore} aggregates unbalanced paired outcomes into per-item scores \citet{DavidScore}. Random-walk and spectral constructions include \textsc{RankCentrality} \citet{NS17}, \textsc{SerialRank} \citet{SerialRank}, \textsc{SyncRank} \citet{SyncRank}, and SVD-based scores \citet{SVD_RS_NRS}. Graph scoring methods include eigenvector centrality \citet{EigenvectorCentrality}, \textsc{PageRank} \citet{PageRank}, and \textsc{SpringRank} \citet{SpringRank}. Spectral and synchronization approaches already note connections between ranking inconsistency and feedback structure \citet{SerialRank,SyncRank}.

\subsubsection{GNN-Based Ranking}
Learning-based methods infer ranking scores from directed comparison graphs. \textsc{GNNRank} uses directed-GNN backbones with ranking-oriented losses \citet{he22}. We compare against the framework's \textsc{DIGRAC} \citet{DIGRAC} and inception-style (\textsc{ib}) \citet{DiGCN} backbones. These methods can be accurate but incur training cost and hyperparameter sensitivity, which motivates a training-free alternative under fixed wall-clock budgets.

\subsubsection{Local-Ratio Cycle Breaking, Add-Back Lineage, and Structural Repair}
Demetrescu and Finocchi \citet{CI03} give a local-ratio cycle-breaking heuristic with an approximation guarantee parameterized by the length of a longest simple directed cycle, building on the local-ratio framework of \citet{BK04}. Their Phase~2 reinserts deleted arcs subject to an \emph{exact} cycle-safety test (equivalently: reinsert $(u,v)$ iff $v$ does not reach $u$ in the current DAG), and they remark that ordering candidates by decreasing weight can improve solution quality. Exact cycle-safe reinsertion and weight-prioritized ordering are therefore prior; they are not newly invented here.

\subsubsection{Relation to \citet{VK25}}
Prior work by \citet{VK25} operationalizes ranking via the \citet{CI03} local-ratio backbone, adopts descending-weight reinsertion under the exact cycle-safety test, and includes order-preserving ternary score refinement. The accompanying practical codebase later used a weaker fixed-topological-position proxy for reinsertion (accepting $(u,v)$ only if it is forward under one fixed topological order). That proxy is a sufficient but not necessary condition for cycle safety and can reject admissible arcs; multipass INS1/INS2/INS3 variants in the submitted manuscript compensate partially for this weakening and are retained only as historical comparison points. This revision restores exact reachability-aware reinsertion as an implementation-fidelity correction, adds formal ranking$\leftrightarrow$MWFAS analysis and guarantee/complexity qualifications, introduces a secondary conflict-triggered weighted min-cut exchange, and expands classical/GNN empirical validation.

\begin{table}[t]
\centering
\small
\setlength{\tabcolsep}{3.5pt}
\begin{tabular}{@{}lccc@{}}
\toprule
\textbf{Component} & \textbf{DF03} & \textbf{VK25} & \textbf{This work} \\
\midrule
Local-ratio backbone & yes & yes & prior (used) \\
Exact cycle-safe reinsertion & yes & yes (spec.) & fidelity restore \\
Weight-ordered reinsertion & suggested & yes & prior (used) \\
Time-bounded practical impl.\ & -- & partial & audited \\
Order-preserving refinement & -- & yes & prior (used) \\
Formal ranking$\leftrightarrow$MWFAS & -- & asserted & proved \\
DF03 guarantee under timeout & N/A & not audited & qualified (no) \\
Weighted min-cut exchange & -- & -- & secondary new \\
Broad classical + GNN eval.\ & -- & GNN-focused & expanded \\
\bottomrule
\end{tabular}
\caption{Novelty separation relative to Demetrescu--Finocchi (DF03) and \citet{VK25}. ``Fidelity restore'' means restoring the established exact test after an implementation weakening, not inventing the test.}
\label{tab:novelty_separation}
\end{table}

\section{Method and Theory}
\label{sec:framework}

We describe a deterministic, training-free ranking pipeline for weighted pairwise-comparison graphs. The canonical revised pipeline is:
\begin{enumerate}
\item \textbf{Phase~A (prior):} local-ratio-style cycle breaking \citet{CI03,VK25}.
\item \textbf{Phase~B (fidelity correction):} exact weight-prioritized reachability-aware reinsertion \citet{CI03,VK25}.
\item \textbf{Optional secondary repair (new):} conflict-triggered weighted min-cut exchange.
\item \textbf{Phase~C (prior lineage):} topological score extraction and optional order-preserving refinement \citet{VK25}.
\end{enumerate}
Legacy multipass topo-proxy variants (INS1/INS2/INS3) appear only as submitted-manuscript comparison points; they are not the revised method's central contribution. We do not claim a new general MWFAS approximation algorithm.

\subsection{Problem Setting}
\label{subsec:framework_problem}
Let $V=\{1,\dots,n\}$ be items and $G=(V,E,w)$ a directed graph with strictly positive arc weights $w:E\to\mathbb{R}_{>0}$. Antiparallel arcs are allowed. The solver input is a sparse weighted adjacency matrix converted to arrays $(\texttt{src},\texttt{dst},\texttt{w})$ of length $m=|E|$. The output is a score vector $s\in\mathbb{R}^n$ with larger values indicating higher rank. Formal ranking and MWFAS definitions appear in \S\ref{sec:intro}.

\subsection{Ranking--MWFAS Optimum-Value Equivalence}
\label{subsec:theory_equivalence}

\begin{proposition}[Optimum-value equivalence]
\label{prop:rank_mwfas}
Let $G=(V,E,w)$ with $w>0$. Write
\[
\mathrm{OPT}_{\mathrm{rank}}
=
\min_{R}\sum_{(i,j)\in E:\,R(i)>R(j)} w_{ij},
\qquad
\mathrm{OPT}_{\mathrm{MWFAS}}
=
\min\bigl\{\textstyle\sum_{e\in F}w_e : F\subseteq E,\ (V,E\setminus F)\ \text{acyclic}\bigr\}.
\]
Then $\mathrm{OPT}_{\mathrm{rank}}=\mathrm{OPT}_{\mathrm{MWFAS}}$.
The correspondence between optimal rankings and optimal feedback arc sets is many-to-many: we do \emph{not} claim a bijection between arbitrary rankings and arbitrary feedback arc sets.
\end{proposition}

\begin{proof}
Write $L(R)=\sum_{(i,j)\in E:\,R(i)>R(j)} w_{ij}$ for the ranking loss of $R$.

($\mathrm{OPT}_{\mathrm{MWFAS}}\le\mathrm{OPT}_{\mathrm{rank}}$.)
Fix any ranking $R$ and let $F_R=\{(i,j)\in E:R(i)>R(j)\}$. Every retained arc is forward under $R$, so $R$ is a topological order of $(V,E\setminus F_R)$ and $F_R$ is a feedback arc set of weight $L(R)$. Hence $\mathrm{OPT}_{\mathrm{MWFAS}}\le L(R)$ for every $R$, and taking the best ranking yields $\mathrm{OPT}_{\mathrm{MWFAS}}\le\mathrm{OPT}_{\mathrm{rank}}$.

($\mathrm{OPT}_{\mathrm{rank}}\le\mathrm{OPT}_{\mathrm{MWFAS}}$.)
Let $F^\star$ be an optimal feedback arc set and let $R$ be any topological order of the DAG $(V,E\setminus F^\star)$. Every retained arc is forward under $R$, so the set of backward arcs of $R$ is contained in $F^\star$. Therefore $L(R)\le w(F^\star)=\mathrm{OPT}_{\mathrm{MWFAS}}$, which implies $\mathrm{OPT}_{\mathrm{rank}}\le\mathrm{OPT}_{\mathrm{MWFAS}}$.
\end{proof}

\begin{remark}
If $F$ is inclusion-minimal, then for every topological order $R$ of $(V,E\setminus F)$ one has $L(R)=w(F)$. Optimal feedback arc sets under strictly positive weights are inclusion-minimal, but the optimum-value claim above does not require uniqueness of $F^\star$ or of $R$.
\end{remark}

\subsection{Unified Pipeline Pseudocode}
\label{subsec:pseudocode}

\begin{algorithm}[H]
\caption{Canonical revised ranking pipeline (training-free)}
\label{alg:ours_canonical}
\KwIn{Weighted digraph $G=(V,E,w)$; wall-clock budget $T$; zero tolerance $\tau$; optional min-cut flag; refinement budgets}
\KwOut{Score vector $s:V\to\mathbb{R}$}
$t_0\leftarrow\mathrm{now}()$;\ $r_e\leftarrow w_e$,\ $\mathrm{alive}_e\leftarrow\mathrm{true}$ for all $e\in E$\;
\While{$\mathrm{true}$}{
  \If{$\mathrm{now}()-t_0>T$}{\textbf{break}\tcp*{may exit while cyclic}}
  $C\leftarrow\mathrm{FindOneCycle}(\mathrm{alive})$ with deterministic vertex/edge order\;
  \If{$C=\emptyset$}{\textbf{break}\tcp*{acyclic}}
  $\Delta\leftarrow\min_{e\in C} r_e$\;
  \eIf{$\Delta\le 0$}{force-kill $\arg\min_{e\in C} r_e$}{
    $r_e\leftarrow r_e-\Delta$ for $e\in C$\;
    delete $\{e\in C:r_e\le\tau\}$;\ \If{none deleted}{force-kill $\arg\min_{e\in C} r_e$}
  }
}
$\mathrm{kept}\leftarrow\mathrm{alive}$\tcp*{Phase~A output}
sort excluded arcs by descending $w$, stable ties by edge id\;
\ForEach{excluded arc $e=(u,v)$ in that order}{
  \If{$\mathrm{now}()-t_0>T$}{\textbf{break}}
  \If{$v\not\rightarrow^\star u$ in current kept DAG}{$\mathrm{kept}[e]\leftarrow\mathrm{true}$}
}
\If{min-cut repair enabled}{apply Alg.~\ref{alg:mincut_exchange} under remaining budget}\;
$\pi\leftarrow\mathrm{TopoSort}(\mathrm{kept})$;\ \If{$\pi=\mathrm{None}$}{$\pi\leftarrow(1,\dots,n)$\tcp*{identity fallback}}
$s_v\leftarrow\max\{1,n-\mathrm{pos}_\pi(v)\}$\;
optionally apply order-preserving refinement under remaining budget\;
\Return{$s$}
\end{algorithm}

\subsection{Phase~A: Local-Ratio Cycle Breaking (Prior)}
\label{subsec:framework_phaseA}
Phase~A follows the local-ratio reduction of \citet{CI03}: maintain residuals, find a directed cycle by deterministic DFS (vertices in increasing index order; outgoing arcs in stored order), subtract the cycle's minimum residual, and delete arcs whose residual falls to at most $\tau$. Forced-progress safeguards kill one minimum-residual arc if floating-point reduction would otherwise stall. Phase~A stops when the alive subgraph is acyclic or when the wall-clock budget expires. Acyclicity is guaranteed only on the former exit.

\subsection{Phase~B: Exact Reachability-Aware Reinsertion (Fidelity Correction)}
\label{subsec:framework_phaseB}
After Phase~A, excluded arcs are scanned once in stable descending original weight. Arc $(u,v)$ is reinserted iff there is no directed path $v\rightarrow^\star u$ in the current kept DAG. This is the exact cycle-safety test of \citet{CI03,VK25}. The submitted fixed-topological-position proxy (accept iff $\mathrm{pos}(u)<\mathrm{pos}(v)$ for one fixed topological order) is weaker and is retained only for historical comparison (legacy INS multipass variants).

\begin{proposition}[Exact add-back correctness]
\label{prop:reachability_addback}
Assume the Phase~A kept graph is a DAG. Consider a single descending-weight scan that inserts $(u,v)$ iff $v\not\rightarrow^\star u$ in the current kept DAG.
\begin{enumerate}
\item \textbf{Acyclicity.} Every accepted insertion preserves acyclicity.
\item \textbf{Monotonicity.} Insertions only enlarge reachability.
\item \textbf{One-pass sufficiency.} An arc rejected when scanned cannot become admissible later in the same monotone scan.
\item \textbf{Single-edge inclusion minimality.} After a complete exact scan, no excluded arc can be restored individually without creating a cycle.
\end{enumerate}
This is inclusion-minimality under single-edge restoration, not a claim of globally minimum-weight FAS.
\end{proposition}

\begin{proof}
(1) If adding $(u,v)$ created a cycle, the remainder of that cycle would be a $v\rightarrow^\star u$ path in the previous DAG, contradicting the acceptance test.
(2) Adding arcs cannot destroy existing paths.
(3) Suppose $(u,v)$ is unsafe at its turn, so a path $v\rightarrow^\star u$ exists. Later insertions only add arcs, so that path persists and $(u,v)$ remains unsafe.
(4) Completeness of the scan means every excluded arc was rejected exactly because restoring it alone would close a cycle.
\end{proof}

\subsection{Optional Secondary Repair: Weighted Min-Cut Exchange}
\label{subsec:framework_mincut}
Some high-weight excluded arcs remain unsafe because $v$ reaches $u$. For such an arc $e=(u,v)$, let $C$ be a minimum-capacity directed $v\rightarrow u$ cut in the retained DAG (capacities $=$ kept-arc weights). The candidate exchange is
\[
D'=(D\setminus C)\cup\{e\}.
\]
Accept only if $w(C)+\mathrm{tol}<w(e)$.

\begin{algorithm}[H]
\caption{Conflict-triggered weighted min-cut exchange (optional)}
\label{alg:mincut_exchange}
\KwIn{Kept DAG $D$; excluded arcs; budgets on candidates/acceptances; tolerance $\mathrm{tol}$}
\ForEach{unsafe candidate $e=(u,v)$ under a fixed deterministic order}{
  $C\leftarrow$ directed min $v\rightarrow u$ cut in $D$\;
  \If{$w(C)+\mathrm{tol}<w(e)$}{
    $D\leftarrow(D\setminus C)\cup\{e\}$\tcp*{strict weighted-FAS improvement}
  }
}
\end{algorithm}

\begin{proposition}[Min-cut exchange]
\label{prop:mincut}
Whenever an exchange is accepted under the strict rule $w(C)+\mathrm{tol}<w(e)$:
\begin{enumerate}
\item $D'$ is acyclic;
\item the removed-FAS weight changes by $\Delta=w(C)-w(e)<0$;
\item repeated accepted exchanges terminate after finitely many steps.
\end{enumerate}
We claim no global optimality and no approximation ratio for this operator. The guarantee is structural improvement of the weighted feedback objective; ranking metrics need not improve monotonically.
\end{proposition}

\begin{proof}
(1) Removing $C$ destroys every $v\rightarrow u$ path, so adding $(u,v)$ cannot close a cycle.
(2) Relative to the previous feedback set, arcs in $C$ re-enter the feedback set and $e$ leaves it, giving $\Delta=w(C)-w(e)$; acceptance forces $\Delta<0$.
(3) There are finitely many subsets of $E$, and each acceptance strictly decreases a nonnegative objective bounded below by $0$.
\end{proof}

\subsection{Phase~C: Scores and Optional Refinement (Prior Lineage)}
\label{subsec:framework_phaseC}
From a kept DAG we compute a deterministic topological order (Kahn's algorithm with index-ordered initial queue) and assign $s_v=\max\{1,n-\mathrm{pos}(v)\}$. If the kept graph is still cyclic (e.g., Phase~A timeout), the implementation falls back to the identity order $(1,\dots,n)$; see \S\ref{subsec:theory_fallback}. Optional adjacent-swap refinement may change the order under a local budget. Optional order-preserving ternary refinement adjusts magnitudes without changing the induced order, following the refinement lineage in \citet{VK25}.

\subsection{Termination, Parameters, and Determinism}
\label{subsec:parameters}
\begin{table}[H]
\centering
\small
\begin{tabular}{@{}llp{6.6cm}@{}}
\toprule
\textbf{Parameter} & \textbf{Role} & \textbf{Notes} \\
\midrule
$T$ (\texttt{time\_limit\_sec}) & Global wall-clock budget & Shared by Phases~A--C; early exit may leave cycles \\
$\tau$ (\texttt{zero\_tol}) & Residual kill threshold & Default $10^{-15}$; numerical, not a modeling knob \\
DFS / Kahn order & Deterministic traversal & Vertices $1..n$; stable edge order \\
Add-back order & Descending $w$, stable ties & Exact reachability test (canonical) \\
Legacy $P\in\{1,2,3\}$ & Submitted topo-proxy passes & Historical INS1/2/3; not canonical \\
Min-cut budgets & Optional repair caps & Secondary; regime-specific benefit \\
Refinement budgets & Phase~C time/passes & Order-preserving ternary from prior lineage \\
Identity fallback & Score extraction if cyclic & No nontrivial general error bound (\S\ref{subsec:theory_fallback}) \\
\bottomrule
\end{tabular}
\caption{Consolidated parameters and termination behavior.}
\label{tab:parameters}
\end{table}

\subsection{DF03 Approximation Guarantee: Idealized vs.\ Practical}
\label{subsec:theory_df03}
\paragraph{Idealized local-ratio Phase~A.}
If Phase~A is executed under the assumptions of the published Demetrescu--Finocchi local-ratio procedure and runs to the relevant completion condition (fully reduced acyclic residual), the published $\lambda$-approximation on removed weight applies to that idealized algorithm, where $\lambda$ is governed by longest-cycle length in the source analysis \citet{CI03}.

\paragraph{Practical implementation.}
The implemented pipeline additionally uses finite wall-clock termination, floating-point zero tolerance, forced-progress safeguards, and a fallback ranking when topological sort fails. Therefore the original theorem does \emph{not} automatically imply that every practical finite-budget output satisfies the same approximation guarantee.

\begin{remark}[No inherited guarantee under premature termination]
\label{rem:no_df03_timeout}
\textbf{We do not claim the DF03 approximation guarantee for prematurely terminated practical runs.}
In particular, a wall-clock cutoff inside Phase~A may leave the kept graph cyclic; subsequent identity-order fallback is outside the theorem's scope.
\end{remark}

\subsection{Fallback Ordering: Unbounded Multiplicative Gap}
\label{subsec:theory_fallback}

\begin{proposition}[No nontrivial multiplicative bound for identity fallback]
\label{prop:fallback}
For every $M>0$ there exists a weighted digraph $G$ with $\mathrm{OPT}_{\mathrm{rank}}>0$ such that if $R_{\mathrm{fb}}$ denotes the identity ranking $R_{\mathrm{fb}}(v)=v$, then
\[
\frac{L(R_{\mathrm{fb}})}{\mathrm{OPT}_{\mathrm{rank}}}>M.
\]
\end{proposition}

\begin{proof}
Fix $\varepsilon\in(0,1)$ and an integer $n\ge 3$ large enough that $\binom{n}{2}/\varepsilon>M$.
On vertex set $\{1,\dots,n\}$, include every arc $(j,i)$ with $j>i$ and weight $1$, and add the single arc $(1,2)$ with weight $\varepsilon$.
Under the identity ranking $R_{\mathrm{fb}}(v)=v$, every arc $(j,i)$ with $j>i$ is backward, so $L(R_{\mathrm{fb}})=\binom{n}{2}$.
Under the reverse ranking $R^\star(v)=n+1-v$, every unit-weight arc is forward and only $(1,2)$ is backward, so $L(R^\star)=\varepsilon$ and thus $\mathrm{OPT}_{\mathrm{rank}}\le\varepsilon$.
Hence $L(R_{\mathrm{fb}})/\mathrm{OPT}_{\mathrm{rank}}\ge\binom{n}{2}/\varepsilon>M$.
(The same instances show that additive regret $L(R_{\mathrm{fb}})-\mathrm{OPT}_{\mathrm{rank}}$ is likewise unbounded as $n$ grows with $\varepsilon$ fixed.)
\end{proof}

\subsection{Complexity}
\label{subsec:framework_complexity}
Let $n=|V|$ and $m=|E|$. For the audited Phase~A implementation, each cycle search scans the static adjacency of all original arcs (dead arcs are skipped but not compacted), costing $O(n+m)$ per iteration, and at most $m$ iterations occur before acyclicity in the unbudgeted case. Therefore the idealized full-completion Phase~A bound is
\[
O(mn+m^2),
\]
which is weaker than a textbook $O(nm)$ accounting that assumes compacted alive adjacency. A wall-clock budget caps completed iterations empirically; it is an execution policy, not an average-case complexity theorem.

Exact reachability add-back costs $O(m\log m)$ for sorting plus, in the dense incremental-closure realization used for moderate $n$, initialization and per-acceptance updates polynomial in $n$ (or BFS queries in a sparse fallback mode). Optional min-cut repair adds per-candidate directed min-cut computations on the kept DAG and is budgeted by candidate/acceptance caps. Phase~C refinement is explicitly time-bounded; worst-case pairwise structures for ratio refinement can reach $O(n^2)$ memory. Empirically, sparse and moderately dense instances finish quickly under the experimental budgets, while the near-complete Finance graph is a large-dense boundary case for Phase~A.

\section{Experimental Results}
\label{sec:exp}
% -------------------------
% -------------------------
\subsection{Experimental Setup}
\label{subsec:exp_setup}

\paragraph{Datasets.}
We evaluate the proposed method on a benchmark suite of $80$ weighted directed comparison graphs drawn from several application domains, following the benchmark setting used in recent ranking-from-comparisons work \citep{he22}. The suite is dominated by temporal sports graphs and includes $60$ basketball instances (seasons 1985--2014 together with a higher-resolution ``finer'' version of each season) and $12$ football instances (seasons 2009--2014 together with finer versions). In addition, the suite contains three faculty hiring networks (Business, Computer Science, and History), one animal social network, two head-to-head competition graphs, one large finance graph, and one additional non-sports instance of ERO style. Across the full suite, the number of vertices ranges from $n=20$ to $n=1315$, and the number of directed edges ranges from $m=107$ to approximately $1.73\times 10^{6}$. Table~\ref{tab:dataset_suite_summary} summarizes the dataset families and their scale ranges.

\begin{table}[t]
\centering
\small
\setlength{\tabcolsep}{4pt}
\renewcommand{\arraystretch}{1.08}
\begin{tabular}{l r r r r r}
\toprule
\textbf{Family} & \textbf{\#} & \multicolumn{2}{c}{\textbf{$n$ range}} & \multicolumn{2}{c}{\textbf{$m$ range}} \\
\cmidrule(lr){3-4}\cmidrule(lr){5-6}
& & \textbf{min} & \textbf{max} & \textbf{min} & \textbf{max} \\
\midrule
Basketball (temporal + finer) & 60 & 282 & 351  & 2{,}904    & 7{,}650 \\
Football (temporal + finer)   & 12 & 20  & 20   & 107        & 380 \\
Faculty (Business)            & 1  & 113 & 113  & 1{,}787    & 1{,}787 \\
Faculty (Computer Science)    & 1  & 206 & 206  & 1{,}407    & 1{,}407 \\
Faculty (History)             & 1  & 145 & 145  & 1{,}204    & 1{,}204 \\
Animal social network         & 1  & 21  & 21   & 193        & 193 \\
Head-to-head competition      & 2  & 602 & 602  & 5{,}010    & 5{,}010 \\
Finance                       & 1  & 1{,}315 & 1{,}315 & 1{,}729{,}225 & 1{,}729{,}225 \\
Other (ERO-style)             & 1  & \multicolumn{2}{c}{within overall range} & \multicolumn{2}{c}{within overall range} \\
\bottomrule
\end{tabular}
\caption{Benchmark suite summary by dataset family. Each instance is a weighted directed graph of pairwise comparisons.}
\label{tab:dataset_suite_summary}
\end{table}

\paragraph{Methods compared.}
We compare four variants of the proposed pipeline: \textsc{OURS-MFAS}, \textsc{OURS-MFAS-INS1}, \textsc{OURS-MFAS-INS2}, and \textsc{OURS-MFAS-INS3}. We also compare against ten classical ranking baselines: \textsc{SpringRank}, \textsc{SyncRank}, \textsc{SerialRank}, \textsc{BTL}, \textsc{DavidScore}, \textsc{EigenvectorCentrality}, \textsc{PageRank}, \textsc{RankCentrality}, \textsc{SVD-RS}, and \textsc{SVD-NRS}. In addition, we include learning-based methods from the GNNRank framework using two directed graph neural network backbones, \textsc{DIGRAC} and \textsc{ib}. For each backbone, we report two fixed benchmark configurations: a distance-based configuration and a proximal-baseline configuration with distance-based pretraining. The main result tables report per-method results explicitly. When we additionally report compact ``best-of'' summaries, we treat them only as oracle-style envelopes over a fixed candidate set and not as standalone methods.

\paragraph{Evaluation metrics.}
We report three upset-based losses computed from the weighted adjacency matrix $A$ and the predicted score vector $s$, where larger scores indicate higher rank. Let $M=A-A^\top$, and define score differences $T_{ij}=s_i-s_j$ over pairs with $M_{ij}\neq 0$. \emph{Upset-naive} is the fraction of pairs for which the predicted sign disagrees with $\mathrm{sign}(M_{ij})$. \emph{Upset-simple} averages the squared disagreement between $\mathrm{sign}(T_{ij})$ and $\mathrm{sign}(M_{ij})$. \emph{Upset-ratio} normalizes both score differences and pairwise edge imbalance by their corresponding sums and averages the resulting squared deviation:
\[
T^{(r)}_{ij}=\frac{s_i-s_j}{s_i+s_j+\varepsilon},
\qquad
M^{(r)}_{ij}=\frac{A_{ij}-A_{ji}}{A_{ij}+A_{ji}+\varepsilon},
\]
where $\varepsilon>0$ is a small constant.

\paragraph{Protocol and compute.}
All methods are evaluated with \texttt{TRIALS=10} repeated runs for each dataset--method pair. For all real datasets in the suite, the dataset seed is fixed internally to \texttt{seeds=[10]}. As a result, deterministic non-GNN methods produce identical outputs across trials apart from timing variability and timeout effects, whereas GNN methods can vary across runs because of stochastic training. Experiments are executed on a SLURM cluster using a typical allocation of one task, $48$ CPU cores, and $100$~GB memory. Non-GNN methods, including classical baselines and the proposed variants, are run under a per-run wall-clock budget of $1800$\,s, while GNN methods are run under a $7200$\,s budget. Timed-out runs are recorded as missing values and are excluded from aggregate statistics via NaN-robust averaging.

In addition to these native budgets, we also report a compute-matched view in which all families are restricted to runs that finish within $1800$\,s. Coverage is reported explicitly in that view. Under the recorded protocol, the compute-matched dataset set is a strict subset of the full suite because the ERO-style instance does not complete within $1800$\,s.

\begin{table}[t]
\centering
\small
\setlength{\tabcolsep}{5pt}
\renewcommand{\arraystretch}{1.10}
\begin{tabular}{p{3.3cm} c p{7.0cm}}
\toprule
\textbf{Item} & \textbf{Value} & \textbf{Notes} \\
\midrule
Trials per dataset--method & 10 & Deterministic non-GNN methods vary only through timing and timeout effects \\
Seed (real datasets) & 10 & Fixed internally to \texttt{seeds=[10]} \\
Non-GNN timeout & 1800\,s & Applied to classical baselines and OURS variants \\
GNN timeout & 7200\,s & Applied per GNN run \\
Compute-matched budget & 1800\,s & Runs are filtered to runtime $\le 1800$\,s and coverage is reported explicitly \\
GNN configurations per backbone & 2 & \texttt{train\_with} in \texttt{\{dist, proximal\_baseline\}} \\
GNN backbones & 2 & \textsc{DIGRAC}, \textsc{ib} \\
OURS metric coverage & 77/80 & Finance times out for all OURS variants under the $1800$\,s non-GNN budget \\
\bottomrule
\end{tabular}
\caption{Protocol and computational settings used in the benchmark suite.}
\label{tab:protocol_summary}
\end{table}

\paragraph{Result aggregation and per-dataset comparisons.}
To ensure numerical consistency in the reported tables, we aggregate accuracy and runtime statistics primarily from saved per-trial arrays whenever they are available, using full-precision means and standard deviations across trials. If a dataset--method entry is unavailable in the array artifacts, we use the corresponding per-dataset summary extracted from the experiment logs and record this fallback in the missingness audit. Unless stated otherwise, all win/tie/loss comparisons are computed from full-precision per-dataset means using a tie tolerance of $|\Delta|\le 10^{-6}$, where $\Delta$ denotes the difference between the losses of two methods on the same dataset. Missing entries, including timeouts and missing runtime values, are reported explicitly through coverage columns and audit tables rather than being omitted silently.

\paragraph{Interpreting mean$\pm$std summaries.}
The manuscript reports two different kinds of variability. When mean$\pm$std summarizes results across datasets, the standard deviation reflects heterogeneity across instances in the benchmark suite. When mean$\pm$std is reported for repeated runs of a fixed dataset--method pair, the standard deviation reflects run-to-run variability. For deterministic non-GNN methods on real datasets, this variability is typically zero except for timing and timeout effects; for GNN methods, it can be nonzero because of stochastic optimization.
% -------------------------
% \subsection{Baselines and Implementation Details}
\label{subsec:baselines_impl}

\paragraph{Unified execution and evaluation.}
All non-GNN methods, including the classical baselines and the proposed variants, are evaluated through the same non-neural benchmark pathway. Each method produces a score vector $s\in\mathbb{R}^n$, where larger values indicate higher rank. When a method naturally returns only an ordering, we convert that ordering to deterministic rank-based scores so that all methods can be evaluated under the same three upset losses defined in \S\ref{subsec:exp_setup}. Learning-based methods are trained and evaluated separately through the GNNRank framework, but their outputs are scored using the same loss implementations for consistency. Because the dataset seed is fixed for all real datasets (Table~\ref{tab:protocol_summary}), non-GNN methods are deterministic under fixed time limits, whereas GNN runs may vary because of stochastic optimization.

\paragraph{Classical ranking baselines.}
Our classical benchmark suite includes ten standard training-free ranking methods: \textsc{SpringRank}, \textsc{SyncRank}, \textsc{SerialRank}, \textsc{BTL}, \textsc{DavidScore}, \textsc{EigenvectorCentrality}, \textsc{PageRank}, \textsc{RankCentrality}, \textsc{SVD-RS}, and \textsc{SVD-NRS}. These methods map the comparison graph directly to a score vector and span statistical, spectral, centrality-based, and physics-inspired ranking families. We use the default settings provided by the corresponding public implementations and do not tune hyperparameters on a per-dataset basis.

\paragraph{Proposed pipeline and variants.}
The proposed method is a combinatorial, training-free pipeline that first constructs an acyclic comparison backbone and then optionally densifies that backbone while preserving acyclicity. Concretely, it (i) removes arcs to eliminate directed cycles using local-ratio-style residual reduction, (ii) reinserts previously removed arcs in stable descending weight order subject to acyclicity, and (iii) extracts a score vector from the resulting DAG and optionally refines score magnitudes under the same upset-based objectives used for evaluation (see \S\ref{sec:framework}). We report four non-GNN variants: \textsc{OURS-MFAS}, \textsc{OURS-MFAS-INS1}, \textsc{OURS-MFAS-INS2}, and \textsc{OURS-MFAS-INS3}. These variants differ only in the number of add-back passes, with INS1, INS2, and INS3 corresponding to $P=1,2,3$, respectively. Unless stated otherwise, we report \textsc{OURS-MFAS-INS3} as \textbf{OURS} in the headline summaries, while INS1 and INS2 serve as lighter variants of the same pipeline.

\paragraph{Complexity and time-bounded operation.}
Let $n$ and $m$ denote the number of vertices and directed edges in the input graph, and let $P\in\{1,2,3\}$ denote the number of add-back passes. Phase~2 sorts edges once by weight, costing $O(m\log m)$, and then performs up to $P$ passes. Each pass computes a topological ordering of the currently kept graph and scans edges in weight order with $O(1)$ acceptance checks, yielding worst-case complexity $O(m\log m + P(n+m))$. Phase~1 repeatedly detects and breaks directed cycles; each cycle-finding call runs in $O(n+m_{\text{alive}})$ over the currently active edges, and at least one edge is deleted whenever a cycle is found, so the number of iterations is at most $m$ in the worst case. In practice, all phases operate under a fixed wall-clock budget (Table~\ref{tab:protocol_summary}), so our emphasis is on empirical runtime under shared budgets rather than on introducing a new approximation guarantee.

\paragraph{GNNRank methods and configurations.}
To represent learning-based alternatives, we include two directed-GNN backbones from the GNNRank framework: \textsc{DIGRAC} and \textsc{ib}. For each backbone, we report two fixed benchmark configurations: a distance-based configuration and a proximal-baseline configuration with distance-based pretraining. This yields four GNN configurations per dataset in total. In the main tables, these configurations are reported explicitly. When we provide compact ``best-of'' comparisons against GNNRank, we treat them only as oracle envelopes over this fixed candidate set and not as standalone methods.

\paragraph{Timeout handling and missingness.}
The benchmark enforces per-run wall-clock limits of $1800$\,s for non-GNN methods and $7200$\,s for GNN methods (Table~\ref{tab:protocol_summary}). Timed-out runs are recorded as missing values and are excluded from aggregate statistics via NaN-robust averaging. Coverage, defined as the number of datasets contributing to a method’s summary statistics, is reported explicitly throughout the results section. Under the $1800$\,s non-GNN budget, all proposed variants time out on the Finance instance, so \textbf{OURS} contributes metrics on $77$ of the $80$ datasets in the full suite. Any remaining missing values, including missing runtime entries or explicit dataset-level exclusions, are reported in the missingness audit rather than omitted silently.

\paragraph{Aggregation and numerical precision.}
To reduce artifacts from log formatting or rounding, our primary table construction uses full-precision per-trial arrays saved from the benchmark runs whenever these are available. If a dataset--method entry is unavailable in the saved arrays, we fall back to the corresponding per-dataset summary extracted from experiment logs and record that fallback in the audit outputs. All tables and figures in \S\ref{subsec:main_results}--\S\ref{subsec:exp_contribution} are generated from the same per-method CSV artifacts, and an automated validation step is used to check internal consistency of coverage and missingness summaries.

\begin{table}[t]
\centering
\footnotesize
\setlength{\tabcolsep}{4pt}
\renewcommand{\arraystretch}{1.08}
\begin{tabular}{p{2.0cm} p{4.0cm} p{8.0cm}}
\toprule
\textbf{Group} & \textbf{Method(s)} & \textbf{Implementation / key settings (summary)} \\
\midrule
OURS &
\textsc{OURS-MFAS}, \textsc{OURS-MFAS-INS1}, \textsc{OURS-MFAS-INS2}, \textsc{OURS-MFAS-INS3} &
Common implementation with Phase~1 local-ratio-style cycle breaking, Phase~2 stable weight-prioritized add-back, and optional score refinement. INS1/INS2/INS3 correspond to $P=1/2/3$ add-back passes; INS3 is the main reported variant. \\
Classical baselines &
\textsc{SpringRank}, \textsc{SyncRank}, \textsc{SerialRank}, \textsc{BTL}, \textsc{DavidScore}, \textsc{EigenvectorCentrality}, \textsc{PageRank}, \textsc{RankCentrality}, \textsc{SVD-RS}, \textsc{SVD-NRS} &
Executed through the common non-neural evaluation pipeline using public implementations and default settings, without per-dataset hyperparameter tuning. \\
GNNRank &
\textsc{DIGRAC}, \textsc{ib} &
Evaluated under four fixed configurations per dataset: two backbones crossed with distance-based and proximal-baseline training settings; the proximal-baseline configuration uses distance-based pretraining. \\
\bottomrule
\end{tabular}
\caption{Compared method families and primary configurations used in the benchmark suite. Main tables report per-method results explicitly; any ``best-of'' summaries over GNN configurations are treated only as oracle envelopes for context.}
\label{tab:methods_summary}
\end{table}

\subsection{Main Results on Benchmark Datasets}
\label{subsec:main_results}

\paragraph{Per-method leaderboard on the full suite.}
Table~\ref{tab:overall_leaderboard_non_gnn} summarizes the main per-method results for the non-GNN methods on the full benchmark suite of $80$ datasets. Each row reports robust aggregate performance across datasets, including the median and mean of the three upset losses together with explicit coverage. Lower values indicate better performance. We report the proposed variants and the classical baselines individually, since these methods have meaningful coverage across the suite and admit direct comparison under the shared non-GNN protocol.

For readability, we do not include long raw GNN configuration strings in the main leaderboard table. Instead, the learning-based methods are compared separately through compute-matched summaries and best-in-suite comparisons in \S\ref{subsec:exp_contribution}. This keeps the main results section focused on interpretable method-level comparisons while preserving the paper’s broader classical-versus-learning-based evaluation story.

\begin{table}[t]
\centering
\footnotesize
\setlength{\tabcolsep}{4pt}
\renewcommand{\arraystretch}{1.08}
\begin{tabular}{l r r r r r r l}
\toprule
\textbf{Method} &
\textbf{Median} &
\textbf{Mean} &
\textbf{Median} &
\textbf{Mean} &
\textbf{Median} &
\textbf{Mean} &
\textbf{Coverage} \\
&
\textbf{Upset-Simple} &
\textbf{Upset-Simple} &
\textbf{Upset-Ratio} &
\textbf{Upset-Ratio} &
\textbf{Upset-Naive} &
\textbf{Upset-Naive} & \\
\midrule
\textsc{OURS-MFAS}        & 0.878049 & 0.887880 & 0.505748 & 0.381891 & 0.219512 & 0.221970 & 77 / 80 \\
\textsc{OURS-MFAS-INS3}   & 0.878049 & 0.887880 & 0.505748 & 0.381891 & 0.219512 & 0.221970 & 77 / 80 \\
\textsc{OURS-MFAS-INS2}   & 0.878049 & 0.887914 & 0.505748 & 0.381874 & 0.219512 & 0.221979 & 77 / 80 \\
\textsc{OURS-MFAS-INS1}   & 0.880927 & 0.888633 & 0.502815 & 0.381915 & 0.220232 & 0.222158 & 77 / 80 \\
\midrule
\textsc{BTL}              & 0.825000 & 0.835385 & 0.631922 & 0.575385 & 0.206250 & 0.208718 & 78 / 80 \\
\textsc{RankCentrality}   & 1.000000 & 1.010641 & 0.646415 & 0.573846 & 0.250000 & 0.252949 & 78 / 80 \\
\textsc{SerialRank}       & 1.075000 & 1.093333 & 0.307267 & 0.296795 & 0.268750 & 0.273333 & 78 / 80 \\
\textsc{SpringRank}       & 1.675000 & 1.679615 & 0.476411 & 0.491154 & 0.418750 & 0.420256 & 78 / 80 \\
\textsc{SyncRank}         & 1.825000 & 1.833462 & 0.476411 & 0.490641 & 0.456250 & 0.461795 & 78 / 80 \\
\textsc{PageRank}         & 1.925000 & 1.934615 & 0.960686 & 0.973974 & 0.481250 & 0.488846 & 78 / 80 \\
\textsc{DavidScore}       & 0.925000 & 0.939103 & 0.632356 & 0.578957 & 0.231250 & 0.234615 & 78 / 80 \\
\textsc{EigenvectorCentrality}
                         & 1.825000 & 1.833077 & 0.960686 & 0.973623 & 0.456250 & 0.461603 & 78 / 80 \\
\textsc{SVD-RS}           & 0.880927 & 0.890513 & 0.790599 & 0.550128 & 0.220232 & 0.222949 & 78 / 80 \\
\textsc{SVD-NRS}          & 0.882925 & 0.898333 & 0.489421 & 0.380128 & 0.220732 & 0.224744 & 78 / 80 \\
\bottomrule
\end{tabular}
\caption{Per-method leaderboard on the full benchmark suite for the proposed variants and classical baselines. Values are medians and means across datasets; lower is better for all losses. Coverage indicates the number of datasets contributing to each row.}
\label{tab:overall_leaderboard_non_gnn}
\end{table}

\paragraph{Compute-matched comparison under a shared $1800$\,s budget.}
To reduce the effect of unequal compute budgets, Table~\ref{tab:compute_matched_leaderboard_non_gnn} reports the same type of summary after restricting attention to runs that complete within the shared non-GNN wall-clock budget of $1800$\,s. The compute-matched set contains $79$ datasets. This view is especially useful for assessing whether observed differences arise from ranking quality or from completion behavior under the fixed budget.

Within this restricted regime, the proposed variants remain competitive with the strongest classical baselines while maintaining explicit coverage accounting. As in the full-suite table, learning-based methods are handled separately through the cross-family comparisons reported later, rather than through raw configuration-heavy leaderboards in the main text.

\begin{table}[t]
\centering
\footnotesize
\setlength{\tabcolsep}{4pt}
\renewcommand{\arraystretch}{1.08}
\begin{tabular}{l r r r r r r l}
\toprule
\textbf{Method} &
\textbf{Median} &
\textbf{Mean} &
\textbf{Median} &
\textbf{Mean} &
\textbf{Median} &
\textbf{Mean} &
\textbf{Coverage} \\
&
\textbf{Upset-Simple} &
\textbf{Upset-Simple} &
\textbf{Upset-Ratio} &
\textbf{Upset-Ratio} &
\textbf{Upset-Naive} &
\textbf{Upset-Naive} & \\
\midrule
\textsc{OURS-MFAS}        & 0.878049 & 0.887880 & 0.505748 & 0.381891 & 0.219512 & 0.221970 & 77 / 79 \\
\textsc{OURS-MFAS-INS3}   & 0.878049 & 0.887880 & 0.505748 & 0.381891 & 0.219512 & 0.221970 & 77 / 79 \\
\textsc{OURS-MFAS-INS2}   & 0.878049 & 0.887914 & 0.505748 & 0.381874 & 0.219512 & 0.221979 & 77 / 79 \\
\textsc{OURS-MFAS-INS1}   & 0.880927 & 0.888633 & 0.502815 & 0.381915 & 0.220232 & 0.222158 & 77 / 79 \\
\midrule
\textsc{SpringRank}       & 0.802724 & 0.811757 & 0.569862 & 0.385207 & 0.200681 & 0.202939 & 78 / 79 \\
\textsc{BTL}              & 0.825000 & 0.835385 & 0.631922 & 0.575385 & 0.206250 & 0.208718 & 78 / 79 \\
\textsc{SVD-RS}           & 0.880927 & 0.890513 & 0.790599 & 0.550128 & 0.220232 & 0.222949 & 78 / 79 \\
\textsc{SVD-NRS}          & 0.882925 & 0.898333 & 0.489421 & 0.380128 & 0.220732 & 0.224744 & 78 / 79 \\
\textsc{RankCentrality}   & 1.000000 & 1.010641 & 0.646415 & 0.573846 & 0.250000 & 0.252949 & 78 / 79 \\
\textsc{SerialRank}       & 1.075000 & 1.093333 & 0.307267 & 0.296795 & 0.268750 & 0.273333 & 78 / 79 \\
\textsc{DavidScore}       & 0.925000 & 0.939103 & 0.632356 & 0.578957 & 0.231250 & 0.234615 & 78 / 79 \\
\textsc{PageRank}         & 1.925000 & 1.934615 & 0.960686 & 0.973974 & 0.481250 & 0.488846 & 78 / 79 \\
\textsc{SyncRank}         & 1.825000 & 1.833462 & 0.476411 & 0.490641 & 0.456250 & 0.461795 & 78 / 79 \\
\bottomrule
\end{tabular}
\caption{Compute-matched per-method leaderboard under a shared $1800$\,s runtime budget. Coverage is computed over the $79$ datasets with at least one recorded run within the budget. Lower is better for all losses.}
\label{tab:compute_matched_leaderboard_non_gnn}
\end{table}

\paragraph{Missingness and timeout summary.}
Coverage differences are an important part of the empirical picture, so Table~\ref{tab:missingness-audit-family} reports a compact dataset-level summary of missingness and timeouts by method family. The main limitation for the proposed method under the shared non-GNN budget is the Finance instance, which does not complete for any of the proposed variants. This explains the gap between the full suite size ($80$) and the OURS coverage ($77$) in the leaderboard tables.

\begin{table}[t]
\centering
\small
\setlength{\tabcolsep}{6pt}
\renewcommand{\arraystretch}{1.10}
\begin{tabular}{l r r r}
\toprule
\textbf{Method family} & \textbf{Datasets with metrics} & \textbf{Datasets with runtime} & \textbf{Datasets with any timeout} \\
\midrule
OURS      & 77 & 77 & 1 \\
classical & 79 & 79 & 2 \\
GNN       & 78 & 78 & 2 \\
other     & 1  & 0  & 1 \\
\bottomrule
\end{tabular}
\caption{Missingness and timeout audit by method family at the dataset level. Counts are computed from the same per-method artifacts used for the leaderboard summaries.}
\label{tab:missingness-audit-family}
\end{table}

\paragraph{Instance-level comparisons and best-of envelopes.}
Aggregate summaries can hide meaningful instance-level variation, so we also examine per-dataset win/tie/loss behavior using full-precision per-dataset means and a tie tolerance of $|\Delta|\le 10^{-6}$. Our primary comparisons remain method-specific. When we additionally report compact ``best-of'' views, we use them only as oracle envelopes over fixed candidate sets, not as standalone algorithms. In particular, later summaries compare the proposed method against a classical oracle envelope and a GNN oracle envelope defined from the fixed candidate sets reported in the benchmark. These summaries provide context for what is achievable within each family while avoiding misleading conflation of multiple methods into a single reported algorithm.

\paragraph{Reproducibility note.}
To verify robustness after recent solver updates, we reran the proposed method on two tight-margin football instances and observed no change in upset-simple or runtime. The reported tables are generated from the same per-method CSV artifacts used throughout the analysis, and internal validation checks were used to confirm consistency of table coverage and missingness summaries.


\subsection{Contribution and Practical Advantages}
\label{subsec:exp_contribution}

This subsection summarizes the practical contribution of the proposed method as a scalable, training-free ranking pipeline under fixed wall-clock budgets. The method constructs a ranking directly from the comparison graph by enforcing acyclicity, selectively restoring high-weight edges while preserving that acyclicity, and then extracting scores consistent with the recovered structure. Accordingly, we do \emph{not} present the method as a new state-of-the-art algorithm for minimum weighted feedback arc set in the sense of new approximation guarantees. Rather, the contribution is a deterministic and deployable ranking pipeline grounded in the MWFAS viewpoint and evaluated under the same upset-based criteria used in recent ranking benchmarks.

Two empirical messages emerge from the benchmark. First, \textsc{OURS-MFAS-INS3} (denoted \textbf{OURS}) is frequently competitive with strong classical ranking baselines on a per-dataset basis. Second, relative to training-based GNNRank configurations, \textbf{OURS} often matches or improves upset losses while providing a clear runtime advantage because it avoids training. Taken together, these findings support the paper’s main practical claim: combinatorial acyclicity-based ranking can offer a favorable accuracy--runtime trade-off on large pairwise-comparison graphs.

\paragraph{Competitiveness relative to classical baselines.}
To place \textbf{OURS} relative to the strongest behavior observed among the classical methods, we compare it against a \emph{best-in-suite} classical reference defined post hoc on each dataset as the minimum upset loss among the ten reported classical baselines. This reference is used only as context and is not treated as a standalone algorithm. The comparison is computed on the $77$ datasets for which \textbf{OURS} returns a valid value under the non-GNN $1800$\,s budget.

Under upset-simple, \textbf{OURS} records $38$ wins and $39$ losses against the best-in-suite classical reference, with median gap $0.0244$ and median relative gap $0.0303$. It lies within $10\%$ of the classical reference on $42/77$ datasets. Under upset-ratio, \textbf{OURS} records $45$ wins and $32$ losses, with median relative gap $-0.0781$, and lies within $10\%$ of the classical reference on $47/77$ datasets. These results indicate that the proposed pipeline is often close to the strongest classical performance observed on a dataset, even though it is not tuned on a per-instance basis.

\paragraph{Accuracy relative to training-based GNNRank configurations.}
We also compare \textbf{OURS} against a \emph{best-in-suite} GNN reference defined post hoc on each dataset as the minimum upset loss among the fixed set of reported GNN configurations. As above, this serves only as compact context and does not replace the per-configuration reporting in the main tables.

Across the same $77$ datasets for which both \textbf{OURS} and the GNN reference are defined, \textbf{OURS} records $45$ wins, $1$ tie, and $31$ losses under upset-simple, with median gap $-0.1822$ and median relative gap $-0.1701$. It is within $10\%$ of the GNN reference on $48/77$ datasets. Under upset-ratio, \textbf{OURS} records $47$ wins and $30$ losses, with median relative gap $-0.2282$, and is within $10\%$ of the GNN reference on $47/77$ datasets. Overall, these comparisons show that a training-free ranking pipeline can remain competitive even against stronger but substantially more expensive learning-based alternatives.

\begin{table}[t]
\centering
\small
\setlength{\tabcolsep}{6pt}
\renewcommand{\arraystretch}{1.10}
\begin{tabular}{l r l l r r l}
\toprule
\textbf{Comparison (best-in-suite reference)} & \textbf{$n$} & \textbf{Metric} & \textbf{W/T/L} & \textbf{Median gap} & \textbf{Median rel.\ gap} & \textbf{Within 10\%} \\
\midrule
OURS vs classical best-in-suite & 77 & upset-simple & 38/0/39 & 0.0244 & 0.0303 & 42/77 \\
OURS vs classical best-in-suite & 77 & upset-ratio  & 45/0/32 & --     & $-0.0781$ & 47/77 \\
\midrule
OURS vs GNN best-in-suite       & 77 & upset-simple & 45/1/31 & $-0.1822$ & $-0.1701$ & 48/77 \\
OURS vs GNN best-in-suite       & 77 & upset-ratio  & 47/0/30 & --     & $-0.2282$ & 47/77 \\
\bottomrule
\end{tabular}
\caption{Summary of per-dataset comparisons using full-precision per-dataset means. ``Best-in-suite'' denotes a post-hoc per-dataset minimum over a fixed candidate set and is used only as context; it does not replace per-method or per-configuration reporting. Ties use $|\Delta|\le 10^{-6}$.}
\label{tab:contribution_summary}
\end{table}

\paragraph{Runtime--accuracy trade-off and deployment advantages.}
The main operational advantage of \textbf{OURS} is that it avoids training and runs under the same non-GNN wall-clock budget as the classical baselines. To quantify this advantage, we compare \textbf{OURS} with the best-in-suite GNN configuration on each dataset for which both runtimes are recorded. On these $76$ datasets, define the speedup
\[
S \;=\; \frac{t_{\text{GNN}}}{t_{\text{OURS}}}\,,
\]
where $t_{\text{GNN}}$ is the runtime of the selected GNN reference and $t_{\text{OURS}}$ is the runtime of \textbf{OURS}.

The resulting speedups are strongly favorable to \textbf{OURS}: the $25$th percentile is $4.72\times$, the median is $10.16\times$, and the $75$th percentile is $18.18\times$, with mean $62.69\times$ due to a heavy tail of very slow GNN runs. Moreover, \textbf{OURS} achieves at least $10\times$ speedup on $38/76$ datasets, at least $50\times$ on $17/76$, and at least $100\times$ on $13/76$ datasets. Using upset-simple as the accuracy criterion, the Pareto breakdown over these $76$ datasets is especially informative: \textbf{OURS} is \emph{better and faster} on $45$ datasets and \emph{worse but faster} on $31$ datasets, and it is never slower than the selected GNN configuration under the recorded protocol. This makes the deployment trade-off clear: the proposed method provides deterministic ranking with substantial runtime savings while still improving accuracy on a substantial fraction of instances.

\begin{table}[t]
\centering
\small
\setlength{\tabcolsep}{7pt}
\renewcommand{\arraystretch}{1.10}
\begin{tabular}{l r}
\toprule
\textbf{Runtime summary (OURS vs GNN best-in-suite)} & \textbf{Value} \\
\midrule
Datasets with both runtimes & 76 \\
Speedup $S=t_{\text{GNN}}/t_{\text{OURS}}$: $P25$ / median / $P75$ & $4.72\times$ / $10.16\times$ / $18.18\times$ \\
Speedup mean & $62.69\times$ \\
Datasets with $S\ge 10\times$ / $50\times$ / $100\times$ & 38 / 17 / 13 \\
Pareto: better \& faster / better but slower & 45 / 0 \\
Pareto: worse but faster / worse \& slower & 31 / 0 \\
\bottomrule
\end{tabular}
\caption{Runtime--accuracy trade-off comparing \textbf{OURS} with the best observed GNN configuration per dataset over datasets with both runtimes recorded. Here the GNN runtime is taken from the best-in-suite post-hoc configuration on each dataset.}
\label{tab:contribution_runtime}
\end{table}

\paragraph{Coverage and limitations.}
The method also has clear limitations. Under the $1800$\,s non-GNN budget, \textbf{OURS} does not complete on the largest extremely dense instances, most notably Finance, and the comparisons above are therefore reported only on datasets for which valid outputs are available. In addition, the dominant accuracy tail relative to both classical and GNN references is concentrated in the higher-resolution \emph{finer basketball} instances. These behaviors are consistent with the design of the pipeline: Phase~1 repeatedly searches for cycles and removes arcs, while Phase~2 performs full edge scans in each pass, and both costs become more demanding as graph density increases. At the same time, finer temporal resolution can amplify local inconsistencies that are difficult to resolve using conservative acyclicity-preserving reinsertion alone. These limitations motivate future improvements such as strongly-connected-component-local refinement, reachability-aware insertion rules, and score-refinement procedures that avoid quadratic pair constructions on dense graphs.
\section{Conclusion and Future Research}
This paper presents a scalable, training-free ranking pipeline for weighted pairwise-comparison graphs based on the connection between ranking inconsistency and the Minimum Weighted Feedback Arc Set (MWFAS) formulation. The proposed method combines a local-ratio-style cycle-breaking stage to construct an acyclic backbone with a stable weight-prioritized add-back procedure (INS1/INS2/INS3 variants) that restores high-weight comparisons while preserving acyclicity. From the resulting graph, the pipeline produces real-valued scores whose induced order is evaluated under the same upset-based losses used in recent ranking benchmarks.

Across a benchmark suite of $80$ instances under fixed wall-clock budgets, the proposed method is competitive with strong classical baselines, including \textsc{BTL}, spectral methods, and centrality-based methods, while also providing a consistent runtime advantage over training-based GNNRank configurations on instances where both sides are recorded. These results support the broader practical message of the paper: combinatorial acyclicity-based structure can serve as an effective, deterministic, and deployable alternative for ranking from pairwise comparisons when training cost, reproducibility, and runtime efficiency are important.

Several directions remain for future research. First, improved add-back and refinement strategies, such as reachability-aware insertion or targeted local search within strongly connected components, may reduce the remaining gap on the higher-resolution ``finer'' instances while preserving the time-bounded design. Second, better scalability on extremely dense graphs will likely require faster cycle detection, incremental maintenance of acyclicity constraints, and refinement procedures that avoid quadratic pair constructions. Third, it would be valuable to study principled tie-handling and uncertainty quantification when the comparison graph admits many valid topological orders or contains weakly connected components. Finally, an important direction is to design ranking objectives and calibration procedures that align more closely with downstream upset-based evaluation while retaining the computational advantages of acyclicity-based ranking.

\section*{Acknowledgments}
Portions of this manuscript and the accompanying code were developed with the assistance of AI-based tools. In particular, ChatGPT (OpenAI) and Gemini were used to help draft and edit text and to support coding, and Cursor and GitHub Copilot were used to assist with software development. The author reviewed, verified, and takes full responsibility for the final manuscript and all reported results.





\section*{Statements and Declarations}

\textbf{Funding}
The author received no financial support for the research, authorship, and/or publication of this article.

\textbf{Competing Interests}
The author declares that he has no competing interests.

\textbf{Ethics Approval}
Not applicable. This article does not report studies involving human participants or animals.

\textbf{Consent to Participate}
Not applicable.

\textbf{Consent for Publication}
Not applicable.

\textbf{Availability of Data and Materials}
The datasets used in the experiments originate from prior benchmarks cited in the manuscript (e.g., \cite{he22}). For reproducibility, the processed versions used in this study are provided as text files containing weighted directed edge lists in the public repository linked below.

\textbf{Code Availability}
All source code, processed data, and scripts required to reproduce the experiments are publicly available at \url{https://github.com/SoroushVahidi/ranking-by-feedback-arc-set}. The repository contains the implementation of the proposed method, the benchmark drivers, and the materials needed to reproduce the reported results. Implementations of the non-proposed baselines and the GNNRank training and evaluation pipeline are based on the public GNNRank codebase of \citet{he22}, with the corresponding adaptations documented in the repository.


\textbf{Author Contributions}
S. Vahidi is the sole author. He conceived the study, developed the algorithms, performed the experiments, and wrote the manuscript.



\bibliographystyle{apalike}  % Use 'apalike' instead of 'spbasic'
\bibliography{references}




\end{document}
